\section{THUẬT TOÁN DECISION TREE}

\subsection{Khái niệm}
Decision Tree (Cây quyết định) là mô hình phân lớp (hoặc hồi quy) dạng cây, trong đó:
\begin{itemize}
    \item Nút trong: tương ứng với một thuộc tính dùng để tách dữ liệu.
    \item Cạnh: tương ứng với giá trị của thuộc tính đó.
    \item Nút lá: tương ứng với nhãn lớp được dự đoán.
\end{itemize}

Ý tưởng chính của thuật toán Decision Tree: Chia để trị. Tại mỗi bước, chọn thuộc tính tốt nhất để tách dữ liệu thành các nhánh con sao cho các mẫu trong cùng một nhánh có cùng nhãn lớp càng nhiều càng tốt.

\subsection{Quy trình phân lớp bằng thuật toán Decision Tree (ID3)}
Quy trình chung:
\begin{enumerate}[label=\textbf{Bước \arabic*:}]
  \item Chuẩn bị và làm sạch dữ liệu.
  \item Tại tập $S$ hiện thời, tính độ bất định (Entropy) và đánh giá các thuộc tính.
  \item Chọn thuộc tính có Information Gain lớn nhất (ID3); với C4.5 sử dụng Gain Ratio.
  \item Tách tập theo thuộc tính được chọn rồi gọi đệ quy cho từng nhánh.
  \item Dừng khi các mẫu đồng nhất nhãn hoặc hết điều kiện chia; cắt tỉa nếu cần.
\end{enumerate}
\begin{align}
H(S)&=-\sum_{c}p(c)\log_2p(c),\\
\operatorname{Gain}(S,A)&=H(S)-\sum_{v\in\operatorname{Values}(A)}\frac{|S_v|}{|S|}H(S_v),\\
\operatorname{Gini}(S)&=1-\sum_c p(c)^2.
\end{align}

\subsection{Mã giả}
\begin{lstlisting}[caption={Mã giả thuật toán ID3}]
ID3(S, Attributes)
Input:
    S: tap mau huan luyen
    Attributes: danh sach thuoc tinh con lai co the tach
Output:
    Mot cay quyet dinh

1. Neu tat ca mau trong S thuoc cung mot lop C:
       Tra ve nut la gan nhan C
2. Neu Attributes rong:
       Tra ve nut la voi nhan = lop xuat hien nhieu nhat trong S
3. Tinh Entropy(S)
4. Voi moi thuoc tinh A trong Attributes:
       Tinh Gain(S, A)
5. Chon A* = thuoc tinh co Gain(S, A) lon nhat
6. Tao nut N tuong ung voi thuoc tinh A*
7. Voi moi gia tri v cua A*:
       S_v = tap con cua S voi A* = v
       Neu S_v rong:
           Tao nut la con cua N, gan nhan = lop da so trong S
       Nguoc lai:
           Gan nhanh con cho N la ket qua goi de quy ID3(S_v, Attributes \ {A*})
8. Tra ve nut N
\end{lstlisting}

\subsection{Ví dụ minh họa - Tính tay chi tiết}

\subsubsection*{Mô tả bài toán}
Bài toán: Dựa vào điều kiện thời tiết, ta quyết định ``Có nên chơi tennis hay không''. \\
Lớp: $\text{PlayTennis} \in \{\text{Yes}, \text{No}\}$. \\
Mục tiêu: từ dữ liệu quá khứ, xây cây quyết định giúp dự đoán chơi / không chơi.

\subsubsection*{Mô tả dữ liệu}
Bảng dữ liệu (14 mẫu kinh điển - PlayTennis):
\begin{table}[H]
    \centering
    \begin{tabular}{|c|c|c|c|c|c|}
        \hline
        \textbf{No} & \textbf{Outlook} & \textbf{Temperature} & \textbf{Humidity} & \textbf{Wind} & \textbf{PlayTennis} \\
        \hline
        1 & Sunny & Hot & High & Weak & No \\ \hline
        2 & Sunny & Hot & High & Strong & No \\ \hline
        3 & Overcast & Hot & High & Weak & Yes \\ \hline
        4 & Rain & Mild & High & Weak & Yes \\ \hline
        5 & Rain & Cool & Normal & Weak & Yes \\ \hline
        6 & Rain & Cool & Normal & Strong & No \\ \hline
        7 & Overcast & Cool & Normal & Strong & Yes \\ \hline
        8 & Sunny & Mild & High & Weak & No \\ \hline
        9 & Sunny & Cool & Normal & Weak & Yes \\ \hline
        10 & Rain & Mild & Normal & Weak & Yes \\ \hline
        11 & Sunny & Mild & Normal & Strong & Yes \\ \hline
        12 & Overcast & Mild & High & Strong & Yes \\ \hline
        13 & Overcast & Hot & Normal & Weak & Yes \\ \hline
        14 & Rain & Mild & High & Strong & No \\ \hline
    \end{tabular}
    \caption{Dữ liệu PlayTennis (Tổng: 14 mẫu -- Yes: 9, No: 5)}
\end{table}

\subsubsection*{Tính tay từng bước theo quy trình ID3}
\textbf{Bước 1.} Entropy toàn bộ tập $S$: Tập $S$ có 9 Yes và 5 No:
$$\text{Entropy}(S) = -\left(\frac{9}{14}\right)\log_2\left(\frac{9}{14}\right) - \left(\frac{5}{14}\right)\log_2\left(\frac{5}{14}\right) \approx 0.940$$

\textbf{Bước 2.} Tính Information Gain cho từng thuộc tính: Ta xét 4 thuộc tính: Outlook, Temperature, Humidity, Wind.
\begin{itemize}
    \item 2.1. Thuộc tính Outlook (Sunny, Overcast, Rain):
    \begin{itemize}
        \item Sunny: (Mẫu 1, 2, 8, 9, 11) $\rightarrow$ 2 Yes, 3 No. $\text{Entropy}(\text{Sunny}) \approx 0.971$
        \item Overcast: (Mẫu 3, 7, 12, 13) $\rightarrow$ 4 Yes, 0 No. $\text{Entropy}(\text{Overcast}) = 0$
        \item Rain: (Mẫu 4, 5, 6, 10, 14) $\rightarrow$ 3 Yes, 2 No. $\text{Entropy}(\text{Rain}) \approx 0.971$
        $$\text{Gain}(S, \text{Outlook}) = 0.940 - \left[\frac{5}{14}(0.971) + \frac{4}{14}(0) + \frac{5}{14}(0.971)\right] \approx 0.246$$
    \end{itemize}
    \item 2.2. Thuộc tính Temperature: $\text{Gain}(S, \text{Temperature}) \approx 0.029$
    \item 2.3. Thuộc tính Humidity: $\text{Gain}(S, \text{Humidity}) \approx 0.152$
    \item 2.4. Thuộc tính Wind: $\text{Gain}(S, \text{Wind}) \approx 0.048$
\end{itemize}
\textbf{Kết luận:} Outlook có Information Gain lớn nhất $\rightarrow$ chọn Outlook làm nút gốc.

\subsubsection*{Xây tiếp các nhánh con}
\begin{itemize}
    \item Nhánh 1: $\text{Outlook} = \text{Overcast}$: Những mẫu có $\text{Outlook} = \text{Overcast}$ đều $\text{PlayTennis} = \text{Yes}$ (4 mẫu). $\text{Entropy} = 0 \rightarrow$ nút lá: Yes.
    \item Nhánh 2: $\text{Outlook} = \text{Sunny}$: Tập con Sunny (Mẫu 1, 2, 8, 9, 11) có 2 Yes, 3 No. $\text{Entropy} \approx 0.971$. Tính Gain cho các thuộc tính còn lại:
    \begin{itemize}
        \item $\text{Gain}(\text{Sunny}, \text{Humidity}) \approx 0.971$ (Lớn nhất) $\rightarrow$ Chọn Humidity làm nút con.
        \item $\text{Humidity} = \text{High}$: (Mẫu 1, 2, 8) $\rightarrow$ Tất cả No $\rightarrow$ Lá: No.
        \item $\text{Humidity} = \text{Normal}$: (Mẫu 9, 11) $\rightarrow$ Tất cả Yes $\rightarrow$ Lá: Yes.
    \end{itemize}
    \item Nhánh 3: $\text{Outlook} = \text{Rain}$: Tập con Rain (Mẫu 4, 5, 6, 10, 14) có 3 Yes, 2 No. $\text{Entropy} \approx 0.971$. Tính Gain:
    \begin{itemize}
        \item $\text{Gain}(\text{Rain}, \text{Wind}) \approx 0.971$ (Lớn nhất) $\rightarrow$ Chọn Wind làm nút con.
        \item $\text{Wind} = \text{Weak}$: (Mẫu 4, 5, 10) $\rightarrow$ Tất cả Yes $\rightarrow$ Lá: Yes.
        \item $\text{Wind} = \text{Strong}$: (Mẫu 6, 14) $\rightarrow$ Tất cả No $\rightarrow$ Lá: No.
    \end{itemize}
\end{itemize}

\subsubsection*{Cây quyết định cuối cùng}
\begin{figure}[H]
\centering
\begin{tikzpicture}[
    node distance=1.5cm and 2.5cm,
    attr/.style={rectangle, rounded corners, draw=black, fill=blue!15, thick, minimum width=2cm, minimum height=0.8cm, align=center},
    leafyes/.style={rectangle, rounded corners, draw=black, fill=green!20, thick, minimum width=1.4cm, minimum height=0.7cm, align=center},
    leafno/.style={rectangle, rounded corners, draw=black, fill=red!20, thick, minimum width=1.4cm, minimum height=0.7cm, align=center},
    edge/.style={->, >=stealth, thick}
]
    \node[attr] (root) {Outlook};
    
    \node[attr, below left=1.4cm and 2.8cm of root] (hum) {Humidity};
    \node[leafyes, below=1.4cm of root] (over) {Yes};
    \node[attr, below right=1.4cm and 2.8cm of root] (wind) {Wind};
    
    \node[leafno, below left=1.3cm and 0.6cm of hum] (hhigh) {No};
    \node[leafyes, below right=1.3cm and 0.6cm of hum] (hnorm) {Yes};
    
    \node[leafyes, below left=1.3cm and 0.6cm of wind] (wweak) {Yes};
    \node[leafno, below right=1.3cm and 0.6cm of wind] (wstrong) {No};
    
    \draw[edge] (root) -- node[above left] {Sunny} (hum);
    \draw[edge] (root) -- node[right] {Overcast} (over);
    \draw[edge] (root) -- node[above right] {Rain} (wind);
    
    \draw[edge] (hum) -- node[above left] {High} (hhigh);
    \draw[edge] (hum) -- node[above right] {Normal} (hnorm);
    
    \draw[edge] (wind) -- node[above left] {Weak} (wweak);
    \draw[edge] (wind) -- node[above right] {Strong} (wstrong);
\end{tikzpicture}
\caption{Cây quyết định ID3 hoàn chỉnh cho bài toán PlayTennis}
\end{figure}

\subsubsection*{Phân tích và kết luận}
\begin{itemize}
    \item Cây trực quan, dễ hiểu, dễ giải thích cho người dùng.
    \item Các nhánh tương ứng với luật IF – THEN:
    \begin{itemize}
        \item IF $\text{Outlook} = \text{Sunny}$ AND $\text{Humidity} = \text{High}$ THEN $\text{Play} = \text{No}$.
        \item IF $\text{Outlook} = \text{Rain}$ AND $\text{Wind} = \text{Weak}$ THEN $\text{Play} = \text{Yes}$.
    \end{itemize}
    \item Trên dữ liệu huấn luyện, cây này phân loại chính xác toàn bộ 14 mẫu.
\end{itemize}

\subsection{Chương trình minh họa Decision Tree}
\lstinputlisting[language=Python,caption={Chương trình Python minh họa thuật toán Decision Tree}]{source_code/decision_tree_85.py}

\subsection{Ưu / Nhược điểm của Decision Tree}
\textbf{Ưu điểm:}
\begin{itemize}
    \item Trực quan, dễ hiểu.
    \item Không cần giả định phân phối dữ liệu.
    \item Xử lý đa dạng loại dữ liệu.
    \item Tốc độ nhanh.
    \item Tự động chọn lọc thuộc tính quan trọng.
\end{itemize}

\textbf{Nhược điểm:}
\begin{itemize}
    \item \textbf{Dễ bị overfitting:} Nếu không cắt tỉa, cây có thể quá sâu và bám vào dữ liệu nhiễu.
    \item \textbf{Không ổn định:} Dữ liệu thay đổi nhỏ có thể làm cây thay đổi cấu trúc lớn.
    \item \textbf{Thiên vị thuộc tính nhiều giá trị:} ID3 thường thiên vị các thuộc tính có nhiều giá trị riêng biệt.
    \item \textbf{Ranh giới quyết định:} Biểu diễn ranh giới phức tạp kém hơn các mô hình như SVM hay Neural Networks.
\end{itemize}